\title{The Era of 1-bit LLMs: All Large Language Models are in 1.58 Bits}

\begin{abstract}
Emerging studies like BitNet~\cite{bitnet} are ushering in the age of 1-bit Large Language Models (LLMs). In this paper, we present a ternary 1-bit LLM version, named \textbf{\text{BitNet b1.58}}, in which every parameter or weight is constrained to the set \{-1, 0, 1\}. This model achieves equivalence with traditional full-precision (FP16 or BF16) Transformer LLMs of identical scale and with the same volume of training data, both in terms of perplexity and downstream task effectiveness. Crucially, BitNet b1.58 offers major advantages in terms of efficiency, including reductions in latency, memory footprint, computational throughput, and overall energy demand. Perhaps most significantly, this 1.58-bit LLM formulation introduces an updated scaling law and provides new methodological guidelines for building future LLMs that maximize both efficiency and performance. Moreover, this innovation enables an alternative computation strategy and facilitates the design of hardware architectures specialized for 1-bit LLMs.
\end{abstract}

\section{The Era of 1-bit LLMs}

The rapid evolution of Large Language Models (LLMs) has substantially expanded the capabilities and scope of artificial intelligence, yielding impressive results across numerous natural language processing benchmarks. However, escalating model sizes create significant obstacles for real-world application, especially due to concerns regarding soaring energy usage and associated economic and environmental implications.

To mitigate these issues, post-training quantization techniques have been developed for producing inference-time low-bit models~\cite{smoothquant, gptq, quip, quip_sharp}. By reducing the numeric precision of both weights and activations, these methods dramatically decrease memory consumption and the computational resources needed for deploying LLMs. The overall movement in the community has shifted from the original 16-bit representations toward even more compact 4-bit formats~\cite{gptq, awq}. Although widely adopted in practice for industrial LLM systems, post-training quantization remains inherently limited and does not yield optimal results.

Recent advancements in 1-bit model design, exemplified by BitNet~\cite{bitnet}, point to a compelling pathway for decreasing the operational costs associated with LLMs without sacrificing their effectiveness. Standard LLM implementations typically utilize 16-bit floating-point representations (such as FP16 or BF16), with matrix multiplication comprising the core computational workload. Consequently, most of the computational expense arises from operations involving floating-point multiplication and addition. In contrast, the architecture of BitNet enables matrix multiplications to be executed solely with integer additions, yielding substantial reductions in energy usage for LLMs. Since the power budget often constrains computational performance in modern chips, these energy reductions directly contribute to faster processing capabilities.

Beyond computation, the transfer of model parameters from off-chip DRAM to the on-chip accelerator’s memory (for example, SRAM) during inference incurs considerable overhead. While increasing SRAM capacity has been explored as a means to boost throughput, this solution is accompanied by cost increments far exceeding those of DRAM. In contrast, 1-bit LLMs require significantly less memory, both in terms of storage and bandwidth, than full-precision models. This pronounced decrease in memory footprint streamlines the process of moving weights from DRAM, yielding lower costs and latency during inference.

In this paper, we present a new 1-bit LLM variant, \textbf{\text{BitNet b1.58}}, characterized by the use of ternary parameters restricted to the set \{-1, 0, 1\}. By incorporating 0 as an additional value beyond the original 1-bit BitNet, this design achieves a representational capacity equivalent to 1.58 bits in binary terms. \text{BitNet b1.58} retains all fundamental advantages of its predecessor, such as a computational paradigm that eliminates most multiplication operations during matrix multiplication, thus facilitating further optimizations. It maintains the same low energy profile as the original 1-bit BitNet and significantly outperforms FP16-based LLMs in memory efficiency, throughput, and latency. Moreover, \text{BitNet b1.58} introduces two noteworthy improvements. Firstly, its representational power is enhanced thanks to explicit feature selection enabled by zero-valued weights, markedly boosting the modeling performance of 1-bit LLMs. Secondly, experimental results demonstrate that \text{BitNet b1.58}, at scales starting from 3B parameters, matches the perplexity and end-task outcomes of full-precision (FP16) baselines under comparable conditions (model size, number of training tokens, etc.).

\section{BitNet b1.58}

\text{BitNet b1.58} builds upon the BitNet framework, a variant of the Transformer model where the standard \emph{nn.Linear} layers are replaced by \emph{BitLinear} modules. This model is trained from the beginning using weights quantized to 1.58 bits and activations at 8 bits. Several alterations distinguish this version from the initial BitNet, which we detail below.

\paragraph{Quantization Function.} To enforce weight values of -1, 0, or +1, we utilize the \emph{absmean} quantization approach. Here, the weight matrix is first normalized by its mean absolute value, followed by rounding each element to its nearest integer in the set \{-1, 0, +1\}:

\begin{equation}
    \widetilde{W} = \mathrm{RoundClip}(\frac{W}{\gamma + \epsilon}, -1, 1),
\end{equation}

\begin{equation}
    \mathrm{RoundClip}(x, a, b) = \max(a, \min(b, \mathrm{round}(x))),
\end{equation}

\begin{equation}
    \gamma = \frac{1}{nm}\sum_{ij} |W_{ij}|.
\end{equation}

For activations, we adopt the quantization method from BitNet, with a key modification: instead of normalizing activations prior to non-linearities within $[0, Q_b]$, we rescale them per token to the interval $[-Q_b, Q_b]$, thereby eliminating the need for zero-point quantization. This adjustment simplifies both implementation and system-level optimization processes and, according to our experiments, has minimal impact on model accuracy.

\paragraph{LLaMA-alike Components.}

With LLaMA~\cite{llama, llama2} serving as the standard architecture for many open-source large language models, we design \text{BitNet b1.58} to be compatible by incorporating components inspired by LLaMA. Our model leverages RMSNorm~\cite{rmsnorm}, SwiGLU~\cite{swiglu}, rotary embeddings~\cite{rope}, and excludes all bias terms. This ensures that \text{BitNet b1.58} can be seamlessly deployed in widely used open-source frameworks (such as Huggingface, vLLM~\cite{vllm}, and llama.cpp\footnote{\url{https://github.com/ggerganov/llama.cpp}}) with minimal modifications.

\section{Results}

We conducted a comparative analysis between \text{BitNet b1.58} and our FP16 implementation of \text{LLaMA LLM} across multiple parameter scales. Both models were pre-trained for 100 billion tokens on the RedPajama dataset~\cite{redpajama} to maintain evaluation parity. Zero-shot task performance was assessed on several benchmarks, including ARC-Easy~\cite{arc}, ARC-Challenge~\cite{arc}, Hellaswag~\cite{hellaswag}, Winogrande~\cite{winoGrande}, PIQA~\cite{piqa}, OpenbookQA~\cite{openbookqa}, and BoolQ~\cite{boolq}. Additionally, validation perplexity was computed on the WikiText2~\cite{wikitext2} and C4~\cite{c4} corpora.

To analyze efficiency, we measured the GPU memory usage and inference latency for both \text{LLaMA LLM} and \text{BitNet b1.58}. These metrics were obtained with the FasterTransformer\footnote{\url{https://github.com/NVIDIA/FasterTransformer}} library, which is specifically optimized for rapid LLM inference on GPUs. For \text{BitNet b1.58}, we utilized the integrated 2-bit kernel from Ladder~\cite{ladder}. The primary focus was on the time required to generate each output token, as this constitutes the principal inference expense.

In Table~\ref{tab:ppl}, we report the perplexity values and associated computational costs for both \text{BitNet b1.58} and \text{LLaMA LLM}. Our results indicate that, at the 3B parameter count, \text{BitNet b1.58} attains perplexity comparable to the full precision \text{LLaMA LLM}, while delivering a 2.71-fold speedup and a 3.55-fold reduction in GPU memory demand. Notably, the 3.9B variant of \text{BitNet b1.58} is 2.4 times faster and uses 3.32 times less memory, all while surpassing the performance of the 3B \text{LLaMA LLM}.

Table~\ref{tab:zero-shot} details the zero-shot accuracy results for various downstream tasks. Evaluation was performed using the \emph{lm-evaluation-harness}\footnote{\url{https://github.com/EleutherAI/lm-evaluation-harness}} framework. Findings show that the performance disparity between \text{BitNet b1.58} and \text{LLaMA LLM} diminishes as model scale increases, with \text{BitNet b1.58} matching full-precision results from the 3B scale upwards. Consistent with the perplexity trends, end-task evaluation demonstrates that the 3.9B version of \text{BitNet b1.58} outperforms \text{LLaMA LLM} 3B, with advantages in both memory usage and latency. Collectively, these results position \text{BitNet b1.58} as a clear Pareto improvement relative to leading LLM baselines.

\paragraph{Memory and Latency}

We extended our experiments by increasing the model sizes to 7B, 13B, and 70B, analyzing the resulting computational expenses. As depicted in Figure~\ref{fig:latency-memory-energy}, both latency and memory utilization display clear trends: performance gains, in terms of reduced latency, are amplified with larger models. Notably, at 70B parameters, \text{BitNet b1.58} demonstrates a 4.1-fold reduction in inference time relative to the \text{LLaMA LLM} baseline. This improvement primarily arises because the runtime associated with \emph{nn.Linear} layers escalates alongside model scale. Memory usage exhibits a comparable pattern, as the embedding layer—which is maintained in full precision—accounts for a decreasing fraction of total memory requirements as models grow. For these measurements, a 2-bit kernel was employed; further cost savings are plausible with additional kernel-level optimizations.

\paragraph{Energy}

In addition, we approximated the energy consumed by arithmetic operations for both \text{BitNet b1.58} and \text{LLaMA LLM}. Our analysis centers on the energy implications of matrix multiplications, the predominant contributor to computational cost in LLMs. Figure~\ref{fig:energy} presents the breakdown of energy usage by operation type. For \text{BitNet b1.58}, most computations involve INT8 additions, whereas \text{LLaMA LLM} relies heavily on FP16 additions and multiplications. Drawing from the energy consumption estimates in~\cite{energycost, pokebnn}, \text{BitNet b1.58} achieves a 71.4-fold reduction in arithmetic operation energy for matrix multiplication on 7nm semiconductor technology. We also assessed end-to-end energy usage across models processing 512 tokens. The findings reveal that as model size increases, the energy efficiency advantage of \text{BitNet b1.58} over the FP16 \text{LLaMA LLM} becomes more pronounced. This is attributed to the escalating contribution of \emph{nn.Linear} operations within larger models, while the influence of other components diminishes with scale.

\paragraph{Throughput}

We assess and contrast the throughput capabilities of \text{BitNet b1.58} and the \text{LLaMA LLM} 70B model using two 80GB A100 GPUs configured with pipeline parallelism~\cite{gpipe}, allowing the large-scale \text{LLaMA LLM} 70B model to fit onto the hardware. The batch size was increased incrementally until reaching the GPU memory ceiling, with all evaluations using a sequence length of 512. As summarized in Table~\ref{tab:throughput}, \text{BitNet b1.58} 70B achieves up to an 11-fold increase in maximum batch size over \text{LLaMA LLM}, which translates to an 8.9x improvement in throughput.

\textbf{\text{BitNet b1.58} introduces a novel scaling law concerning the relationship between model size, inference efficiency, and performance.} By referencing Figures~\ref{fig:latency-memory-energy} and \ref{fig:energy}, we can establish comparative equivalences between 1.58-bit and 16-bit models:

\begin{itemize}
    \item BitNet b1.58 at 13B parameters surpasses 3B FP16 LLMs in latency, memory consumption, and energy requirements.
    \item BitNet b1.58 at 30B parameters demonstrates greater efficiency across latency, memory, and energy usage than 7B FP16 LLMs.
    \item BitNet b1.58 at 70B parameters outperforms 13B FP16 LLMs regarding latency, memory, and energy efficiency.
\end{itemize}

\paragraph{Training with 2T Tokens}

Training token count plays a significant role in large language model effectiveness. To examine the token scalability of \text{BitNet b1.58}, we trained this architecture with 2T tokens, adhering to the data preparation methodology described for StableLM-3B~\cite{StableLM-3B-4E1T}, a leading 3B open-source model. The two models were tested on an evaluation suite including Winogrande~\cite{winoGrande}, PIQA~\cite{piqa}, SciQ~\cite{sciq}, LAMBADA~\cite{lambada}, and ARC-easy~\cite{arc}. Table~\ref{tab:2t} reports the observed zero-shot accuracy; for tasks featuring both accuracy and normalized accuracy, we present the average value. StableLM 3B results at 2T tokens are cited from its corresponding technical report. The evaluation reveals that \text{BitNet b1.58} consistently delivers superior results across all considered benchmarks, highlighting its robust generalization abilities even at the 1.58-bit quantization level.

\section{Discussion and Future Work}

\textbf{1-bit Mixture-of-Experts (MoE) LLMs}

Mixture-of-Experts (MoE) architectures have emerged as an efficient solution for scaling LLMs, offering a substantial reduction in computational FLOPs. Nevertheless, their widespread use is constrained by significant memory requirements and communication costs between processing units. Employing 1.58-bit LLMs provides an effective strategy to overcome these hurdles. Notably, the diminished memory footprint enables deployment across fewer devices, while it also leads to a marked decrease in the costs associated with transferring activation data across networks. If the reduced resource demand allows for model placement entirely within a single device, inter-chip communication becomes a non-issue.

\textbf{Native Support of Long Sequence in LLMs}

The demand for processing extended input sequences has grown with the advancement of LLMs, yet the high memory usage of key-value (KV) caches remains a central bottleneck for long sequence inference. \text{BitNet b1.58} advances native long-context capabilities by compressing activation storage from 16 bits to 8 bits, effectively doubling feasible context length without additional resource requirements. Furthermore, activations could potentially be losslessly encoded in as few as 4 bits or less in 1.58-bit LLMs, a possibility we intend to investigate further.

\textbf{LLMs on Edge and Mobile}

Adopting 1.58-bit LLMs can substantially broaden the applicability of large language models on edge and mobile platforms, which typically suffer from stringent memory and compute constraints. By lowering both memory and energy demands, these models facilitate on-device inference and empower a new suite of applications that had previously been unattainable. This development significantly bolsters the potential of edge and mobile devices for language-related tasks. In addition, 1.58-bit LLMs align well with CPU-based processing—the primary architecture in these environments—enabling efficient execution of \text{BitNet b1.58} and further extending device capabilities.

\textbf{New Hardware for 1-bit LLMs}

Work such as Groq\footnote{\url{https://groq.com/}} highlights the feasibility and promise of specialized hardware, like LPUs, for efficient LLM execution. Building upon these advances, we advocate for targeted hardware and system designs optimized specifically for 1-bit LLMs, taking full advantage of the computational paradigms introduced by BitNet~\cite{bitnet}.